\documentclass[english, parskip=full]{scrartcl}
\usepackage[english]{babel}  % Sprache auf Deutsch 
\usepackage[utf8]{inputenc}  % Kodierung der Datei
\usepackage[left=2cm, right=2cm, top=2cm, bottom=3cm, footskip=1.5cm]{geometry}
\input{myscrarticle}

\newcommand*{\titel}{Lorentz force from electromagnetic field momentum}
\newcommand*{\autor}{Joachim Schwardt}

\hypersetup{pdfauthor={\autor}, pdftitle={\titel}}

%\titlehead{titlehead}
\subject{Theoretical Electrodynamics}
\title{\titel}
\author{\autor}
\date{\today}

\begin{document}

%\begin{titlepage}
\maketitle  % Titel setzen
%\tableofcontents  % Inhaltsverzeichnis setzen
%\end{titlepage}

Consider a point-charge $q$ moving a distance short $\Delta\textbf{r}=\textbf{v}\Delta t$ where $\textbf{v}$ is the velocity and $\Delta t$ is the time-span.
Then the four-potential is given by\footnote{We have set $c=1$.}
\begin{align}
\phi(\textbf{r},t) &= \frac{\mu_0q}{4\pi} \frac{1}{\abs*{\textbf{r} - \textbf{v}t}}, &&& \textbf{A}(\textbf{r},t) &= \textbf{v} \phi(\textbf{r},t).
\end{align}
The electric field is thus given by
\begin{align}
\textbf{E} &= -\nabla \phi - \partial_t \textbf{A} = - \frac{\mu_0 q}{4\pi} \kla*{\nabla + \textbf{v} \partial_t} \frac{1}{\abs*{\textbf{r} - \textbf{v} t}} = \\
&= - \frac{\mu_0 q}{4\pi} \kla*{ -\frac{\textbf{r} - \textbf{v}t}{\abs*{\textbf{r} - \textbf{v} t}^3} - \textbf{v} \frac{\textbf{r} \cdot \textbf{v} - v^2t}{\abs*{\textbf{r} - \textbf{v} t}^3}} = \\
&= \frac{\mu_0q}{4\pi} \frac{1}{\abs*{\textbf{r} - \textbf{v} t}^3} \kla*{ \textbf{r} - \textbf{v} t(1 - v^2) - \textbf{v} (\textbf{r} \cdot \textbf{v})}.
\end{align}
%%%%Assuming an external homogeneous magnetic field $\textbf{B}$, the field momentum becomes
%%%%\begin{align}
%%%%\textbf{P} &= \epsilon_0 \intrnd{3}{\textbf{E} \times \textbf{B}}{r} = -\epsilon_0 \textbf{B} \times \intrnd{3}{ \frac{\mu_0q}{4\pi} \frac{1}{\abs*{\textbf{r} - \textbf{v} t}^3} \kla*{ \textbf{r} - \textbf{v} t(1 - v^2) - \textbf{v} (\textbf{r} \cdot \textbf{v})} } {r}.
%%%%\end{align}
%%%%For $t=0$ this simplifies to
%%%%\begin{align}
%%%%\textbf{P}(t=0) &= -\epsilon_0 \textbf{B} \times \frac{\mu_0q}{4\pi} \intrnd{3}{ \frac{1}{r^3}}{•}
%%%%\end{align}
%%Assuming an external homogeneous magnetic field $\textbf{B}$, the field momentum becomes
%%\begin{align}
%%\textbf{P} &= \epsilon_0 \intrnd{3}{\textbf{E} \times \textbf{B}}{r}.
%%\end{align}
%%Thus the change in the field momentum is determined by the change in the electric field.
%%We may calculate this as
%%\begin{align}
%%\Delta \textbf{E} &= \textbf{E}(t=\Delta t) - \textbf{E}(t=0) = \\
%%&= \frac{\mu_0q}{4\pi} \klb*{\frac{1}{r^3} }• 
%%\end{align}
Assuming an external homogeneous magnetic field $\textbf{B}$, the field momentum becomes
\begin{align}
\textbf{P} &= \epsilon_0 \intrnd{3}{\textbf{E} \times \textbf{B}}{r} = -\epsilon_0 \textbf{B} \times \intrnd{3}{ \frac{\mu_0q}{4\pi} \frac{1}{\abs*{\textbf{r} - \textbf{v} t}^3} \kla*{ \textbf{r} - \textbf{v} t(1 - v^2) - \textbf{v} (\textbf{r} \cdot \textbf{v})} } {r}.
\end{align}
Making the substitution $\textbf{r} \rightarrow \textbf{r} + \textbf{v}t$ yields
\begin{align}
\textbf{P} &= -\frac{q\textbf{B}}{4\pi} \times \intrnd{3}{\frac{\textbf{r} + \textbf{v}^3t - \textbf{v} \kla*{ \textbf{r} \cdot \textbf{v} + v^2t}}{r^3} }{r} = \\
&= -\frac{q\textbf{B}}{4\pi} \times \intrnd{3}{\frac{\textbf{r} - \textbf{v} \kla*{ \textbf{r} \cdot \textbf{v} }}{r^3} }{r}.
\end{align}
The first term vanishes due to spherical symmetry.
The remaining term can be written as
\begin{align}
\textbf{P} &= \frac{q\textbf{B}}{4\pi} \textbf{v} \intrnd{3}{ \frac{\textbf{r} \cdot \textbf{v}}{r^3}  }{r} = \frac{q\textbf{B}}{4\pi} \textbf{v} \kla*{\intrnd{3}{ \frac{\textbf{r}}{r^3}  }{r} \cdot \textbf{v}}
\end{align}
and thus vanishes for the same reason.
Therefore the field momentum is not only time-independent, but actually exactly zero.
Note that this is only the ``interactive'' field momentum, if one were to use the magnetic field generated by the moving charg,e one would get a more complicated result.


%\begin{thebibliography}{9}
%\bibitem{example} description, \url{https://www.exmapleurl}, day.\,month~2021.
%\end{thebibliography}

\end{document}